\documentclass{ctexart}
\usepackage{amsmath, amssymb}
\usepackage{geometry}
\geometry{a4paper, margin=2cm}

\title{\textbf{第11章 变化的磁场和变化的电场} --- 例题}
\author{}
\date{}

\begin{document}
\maketitle

\begin{enumerate}

\item 匀强磁场中，导线可在导轨上滑动，求回路中感应电动势。

\item（书中例题10.2，p.443）一半径 $r = 0.20\,\text{m}$ 的半圆导线和直导线组成一回路，磁场垂直纸面向外，磁感应强度大小 $B = 4t^{2} + 2t + 3$，回路电阻 $R = 2\,\Omega$，其中接一电动势 $\varepsilon = 2.0\,\text{V}$ 的理想电源。求 $t = 10\,\text{s}$ 时回路中的感应电动势的大小和方向及回路中的电流。

\item 在无限长直载流导线的磁场中，有一运动的导体线框，导体线框与载流导线共面，求线框中的感应电动势。

\item（书中例题10.3，p.444）长直导线载有变化的电流 $i = (6t^{2} + 6t)\times10^{-1}\,\text{A}$，接有电源的矩形线框放置在导线的同一平面里。$a = 0.1\,\text{m}$，$b = 0.3\,\text{m}$，$L = 0.3\,\text{m}$，电源电动势 $\varepsilon = 2\,\text{V}$，磁导率 $\mu = 1.46\times10^{-4}\,\text{Tm/A}$，线圈总匝数为 $1000$，电阻 $R = 2\,\Omega$。求 $t = 1\,\text{min}$ 时线圈中电流 $I$ 的大小和方向。

\item 在匀强磁场 $B$ 中，长 $R$ 的铜棒绕其一端 $O$ 在垂直于 $B$ 的平面内转动，角速度为 $\omega$。求棒上的电动势。

\item（书中例题10.5，p.541）圆盘发电机是一个在磁场中转动的金属盘。设圆盘半径 $R = 0.20\,\text{m}$，匀强磁场的磁感应强度 $B = 0.70\,\text{T}$，转速为 $50\,\text{s}^{-1}$。求盘心与盘边缘之间的电势差 $\Delta U$。

\item（书中例题10.6，p.452）磁感应强度为 $B$ 的匀强磁场中，放置一圆形线圈，线圈电阻为 $R$，半径为 $r$，绕直径 $OO'$ 以匀角速度 $\omega$ 旋转，当线圈平面转至与 $B$ 平行时，求 $OA$ 两点的动生电动势及回路中的感应电流。

\item（书中例题10.7，p.453）一通有恒定电流 $I$ 的长直导线，旁边有一个与它共面的三角形线圈 $ACD$，$AC$ 的长为 $l$，$D$ 到 $AC$ 边的垂直距离为 $d$，时刻 $t$，边 $AC$ 与长直导线平行且相距 $r$。试求当线圈以速度 $v$ 沿竖直方向向上运动时，三角形线圈每边上的动生电动势的大小和方向。

\item（书中例题10.8，p.455）半径为 $R$ 的长直螺线管中载有变化电流，管内产生均匀磁场，当磁感应强度以恒定速率 $\dfrac{\partial B}{\partial t}$ 增加时。求：(1) 管内外有旋电场 $E_{\text{旋}}$；(2) 导体棒 $ab$ 两端的电动势。

\item（书中例题10.9，p.457）薄壁导体柱壳，半径为 $r$，高为 $h$，电阻为 $R$，放在 $N$ 匝、长 $L$ 的螺线管中部，$L \gg h$，螺线管中通以交变电流 $I = I_0\sin\omega t$。求柱壳中的感应电流。

\item（书中例题10.11，p.465）空心单层密绕长直螺线管，总匝数为 $N$，长为 $L$，半径为 $R$，且 $L \gg R$。求螺线管的自感。

\item 设一载流回路由两根平行的长直导线组成，求这一对导线单位长度的自感 $L$。

\item（书中例题10.13，p.466）横截面为矩形的密绕螺绕环，总匝数为 $N$，内外半径分别为 $R_1$ 和 $R_2$。求螺绕环的自感。

\item 同轴电缆由半径分别为 $R_1$ 和 $R_2$ 的两个无限长同轴导体柱面组成，求无限长同轴电缆单位长度上的自感。

\item（书中例题10.14，p.468）两个同轴螺线管1和2同绕在一个半径为 $R$ 的长磁介质棒上，绕向相同，截面积等于磁介质棒的截面积，螺线管长分别为 $l_1$ 和 $l_2$，单位长度上的匝数分别为 $n_1$ 和 $n_2$，且 $l_1 \gg R$，$l_2 \gg R$。证明 $M_{21} = M_{12} = M$。

\item（书中例题10.15，p.469）矩形线圈 $ABCD$，长为 $l$，宽为 $a$，匝数为 $N$，放在一长直导线旁边与之共面，矩形线圈中通有电流 $i = I_0\cos\omega t$。求矩形线圈对直导线的互感电动势。

\item（书中例题10.16，p.470）半径分别为 $R$ 和 $r$（$R > r$）的两个同轴线圈，相距为 $d$，且 $d \gg R$，大线圈中通有电流 $I = I_0\sin\omega t$。求：(1) 两线圈的互感系数；(2) 小线圈中的互感电动势。

\item 设平行板电容器极板为圆板，半径为 $R$，两极板间距为 $d$，用缓变电流 $I_C$ 对电容器充电。求 $P_1$、$P_2$ 点处的磁感应强度。

\end{enumerate}

\end{document}
